\chapter{第一章作业}

\section{作业1.1 方程的导出与定解条件}
P10三题
\begin{homework}
[1]. 细杆（或弹簧）受某种外界原因而产生纵振动，以 $u(x,t)$ 表示静止时在 $x$ 点处的点在时刻 $t$ 离开原来位置的偏移。假设振动过程中所产生的张力服从胡克定律，试证明 $u(x,t)$ 满足方程
\[
\frac{\partial}{\partial t}\left(\rho(x)\frac{\partial u}{\partial t}\right)
= \frac{\partial}{\partial x}\left(E\frac{\partial u}{\partial x}\right),
\]
其中 $\rho$ 为杆的密度，$E$ 为杨氏模量。
\end{homework}

\begin{homework}
[6] 绝对柔软而均匀的弦线有一端固定，在它本身重力的作用下，此线处于铅垂的平衡位置，试导出此线的微小横振动方程。
\end{homework}

\begin{homework}
[7] 一柔软均匀的细弦，一端固定，另一端是弹性支承。设该弦在阻力与速度成正比的介质中作微小的横振动，试写出弦的位移所满足的定解问题。
\end{homework}

\section{作业1.2 达朗贝尔公式、波的传播}
P19:5题
\begin{homework}[1]
证明方程
\[
\frac{\partial}{\partial x}\!\left[\Bigl(1-\frac{x}{h}\Bigr)^{2}\frac{\partial u}{\partial x}\right]
= \frac{1}{a^{2}}\Bigl(1-\frac{x}{h}\Bigr)^{2}\frac{\partial^{2}u}{\partial t^{2}}
\]
的通解可以写成
\[
u = \frac{F(x-at)+G(x+at)}{h-x},
\]
并求其满足初始条件
\[
u(x,0)=\varphi(x),\qquad u_t(x,0)=\psi(x)
\]
的解。
\end{homework}

\begin{homework}
[3] 利用传播波法，求解波动方程的古尔萨（Goursat）问题
\[
\begin{cases}
u_{tt}=a^{2}u_{xx}, & -at<x<at,\ t>0,\\[4pt]
u|_{x-at=0}=\varphi(x),\\[4pt]
u|_{x+at=0}=\psi(x),\qquad \varphi(0)=\psi(0).
\end{cases}
\]
\end{homework}

\begin{homework}
[5] 求解
\[
\begin{cases}
u_{tt}-a^{2}u_{xx}=0,\ x>0,\ t>0,\\[4pt]
u(x,0)=\varphi(x),\\[4pt]
u_t(x,0)=0,\\[4pt]
(u_x-ku_t)|_{x=0}=0,
\end{cases}
\]
其中 $k>0$。
\end{homework}

\begin{homework}
[7] 求边值问题
\[
\begin{cases}
u_{tt}-u_{xx}=0,\quad f(t)<x<t,\ t>0,\\[4pt]
u|_{x=t}=\varphi(t),\\[4pt]
u|_{x=f(t)}=\psi(t),
\end{cases}
\]
其中 $\varphi(0)=\psi(0)$，且 $x=f(t)$ 是满足 $|f'(t)|\neq1$ 的光滑曲线。
\end{homework}

\begin{homework}
[8] 求解初值问题
\[
\begin{cases}
u_{tt}-u_{xx}=t\sin x,\qquad t>0,\ -\infty<x<\infty,\\[4pt]
u(x,0)=0,\\[4pt]
u_t(x,0)=\sin x.
\end{cases}
\]
\end{homework}

\section{分离变量法}
P31:4
\begin{homework}[1(2)]
求解初边值问题（$0<x<l,\ t>0$）：
\[
\begin{cases}
u_{tt}=a^{2}u_{xx},\\[3pt]
u(0,t)=0,\quad u_x(l,t)=0,\\[3pt]
u(x,0)=\dfrac{h}{l}x,\quad u_t(x,0)=0.
\end{cases}
\]
\end{homework}

\begin{homework}[2]
求解初边值问题：
\[
\begin{cases}
u_{tt}=a^{2}u_{xx},\\[3pt]
u(0,t)=0,\quad u(l,t)=A\sin^{2}(\omega t),\\[3pt]
u(x,0)=0,\quad u_t(x,0)=0.
\end{cases}
\]
\end{homework}

\begin{homework}
[4] 求解
\[
\begin{cases}
u_{tt}-a^{2}u_{xx}=g,\\[3pt]
u(0,t)=0,\quad u_x(l,t)=0,\\[3pt]
u(x,0)=0,\quad u_t(x,0)=\sin\frac{\pi x}{2l}.
\end{cases}
\]
\end{homework}

\begin{homework}
[6] 求解阻尼波方程
\[
\begin{cases}
u_{tt}+2bu_t=a^{2}u_{xx},\quad b>0,\\[3pt]
u(0,t)=u(l,t)=0,\\[3pt]
u(x,0)=\frac{h}{l}x,\quad u_t(x,0)=0.
\end{cases}
\]
\end{homework}

\section{4 高维柯西问题}
P64：4
\begin{homework}
[1] 利用泊松公式求解三维波动方程的柯西问题：
\[
\begin{cases}
u_{tt}=a^{2}(u_{xx}+u_{yy}+u_{zz}),\\
u(x,y,z,0)=0,\quad u_t(x,y,z,0)=x^{2}+yz.
\end{cases}
\]
以及
\[
\begin{cases}
u_{tt}=a^{2}(u_{xx}+u_{yy}+u_{zz}),\\
u(x,y,z,0)=x^{3}+y^{2}z,\quad u_t(x,y,z,0)=0.
\end{cases}
\]
\end{homework}

\begin{homework}[2]
用降维法导出二维波动方程
\[
u_{tt}=a^{2}(u_{xx}+u_{yy})+f(x,y,t)
\]
在齐次初始条件
\[
u(x,y,0)=0,\qquad u_t(x,y,0)=0
\]
下的求解公式。
\end{homework}

\begin{homework}[5]
求解二维 Klein\text{--}Gordon 型柯西问题：
\[
\begin{cases}
u_{tt}=a^{2}(u_{xx}+u_{yy})+c^{2}u,\\
u(x,y,0)=\varphi(x,y),\\
u_t(x,y,0)=\psi(x,y).
\end{cases}
\]
提示：令 $v=e^{(c/a)z}u$ 并利用三维波动方程。
\end{homework}

\begin{homework}[6]
试用齐次化原理导出非齐次二维波动方程
\[
u_{tt}=a^{2}(u_{xx}+u_{yy})+f(x,y,t)
\]
在初始条件 $u(x,y,0)=0,\ u_t(x,y,0)=0$ 下的求解公式。
\end{homework}

\section{5}
P53
\begin{homework}[1]
试说明：对一维波动方程所描述的波的传播过程一般具有后效现象。
\end{homework}

\begin{homework}[2]
试说明：对一维波动方程，即使初始资料具有紧支集，当 $t\to\infty$ 时其柯西问题的解没有衰减性。
\end{homework}

\begin{homework}[3]
设 $u$ 为初始资料 $\varphi$ 及 $\psi$ 具有紧支集的二维波动方程的解。试证明：对任意固定的 $(x_0,y_0)\in\mathbb{R}^2$，$u(x_0,y_0,t)$ 以 $t^{-1/2}$ 的速度趋于零。
\end{homework}

\section{6}
P62
\begin{homework}[1]
对受摩擦力作用且具有固定端点的有界弦振动，满足方程
\[
u_{tt}=a^{2}u_{xx}-c u_t,
\]
其中常数 $c>0$，证明其能量是减少的，并由此证明方程
\[
u_{tt}=a^{2}u_{xx}-c u_t+f
\]
的初边值问题解的唯一性。
\end{homework}

\begin{homework}[5]
考虑波动方程的第三类初边值问题
\[
\begin{cases}
u_{tt}-a^{2}(u_{xx}+u_{yy})=0, & t>0,\ (x,y)\in\Omega,\\[4pt]
u|_{t=0}=\varphi(x,y),\qquad u_t|_{t=0}=\psi(x,y),\\[4pt]
\left.\left(\dfrac{\partial u}{\partial n}+\sigma u\right)\right|_{\Gamma}=0,
\end{cases}
\]
其中 $\sigma>0$ 为常数，$\Gamma$ 为 $\Omega$ 的边界，$n$ 为 $\Gamma$ 上的单位外法线向量。
对上述定解问题的解，定义能量函数
\[
E(t)=\iint_{\Omega}[u_t^2 + a^{2}(u_x^2+u_y^2)]\,dxdy
\ +\ a^{2}\int_{\Gamma}\sigma u^{2}\,ds.
\]
试证明 $E(t)\equiv\text{常数}$，并由此证明上述定解问题的唯一性。
\end{homework}

\begin{homework}[7]
设 $u(x,t)$ 是 $[0,1]\times[0,\infty)$ 中初边值问题
\[
\begin{cases}
u_{tt}-u_{xx}=0,\\[4pt]
u|_{x=0}=u|_{x=1}=0,\\[4pt]
u|_{t=0}=0,\qquad u_t|_{t=0}=x^{2}(1-x),
\end{cases}
\]
的解。求
\[
\int_{0}^{1}\left[u_t^{2}(x,t)+u_x^{2}(x,t)\right]\,dx.
\]
\end{homework}
\chapter{第二章作业}
\section{2 初边值的分离变量法}[一题]
\begin{homework}[1]
用分离变量法求下列定解问题的解：
\[
\begin{cases}
\dfrac{\partial u}{\partial t}
= a^{2}\dfrac{\partial^{2}u}{\partial x^{2}},
& t>0,\ 0<x<\pi,\\[6pt]
u(0,t)=0,\quad
\dfrac{\partial u}{\partial x}(\pi,t)=0,
& t>0,\\[6pt]
u(x,0)=f(x),
& 0<x<\pi.
\end{cases}
\]
\end{homework}
\begin{homework}[2]
用分离变量法求解热传导方程的初边值问题：
\[
\begin{cases}
\dfrac{\partial u}{\partial t}
= \dfrac{\partial^2 u}{\partial x^2}, & (t>0,\,0<x<1),\\[6pt]
u(x,0)=
\begin{cases}
x, & 0<x\le\dfrac{1}{2},\\[4pt]
1-x, & \dfrac{1}{2}<x<1,
\end{cases}\\[10pt]
u(0,t)=u(1,t)=0, & (t>0).
\end{cases}
\]
\end{homework}


\section{3柯西问题——三题}
% 约定：对 $f\in L^1(\mathbb R)$ 定义一维傅里叶变换
% \[
%   \mathcal F[f](\omega)=\hat f(\omega)
%   =\int_{-\infty}^{\infty}f(x)e^{-i\omega x}\,dx .
% \]

%------------------ 题 1 ------------------%
\begin{homework}[1]
求下列函数的傅里叶变换：
\begin{enumerate}[(1)]
  \item $e^{-\eta x^2}\quad(\eta>0)$；
  \item $e^{-a|x|}\quad(a>0)$；
  \item $\displaystyle\frac{x}{(a^2+x^2)^k},\qquad
        \frac1{(a^2+x^2)^k}\quad(a>0,\;k\in\mathbb N)$。
\end{enumerate}
\end{homework}
\begin{dxtips}
    指数是相加不是相乘。
\end{dxtips}
%------------------ 题 2 ------------------%
\begin{homework}[2]
设 $f(x)$ 在 $(-\infty,\infty)$ 绝对可积，证明其傅里叶变换
\[
\hat f(\omega)=\int_{-\infty}^{\infty}f(x)e^{-i\omega x}\,dx
\]
是连续函数。
\end{homework}



%------------------ 题 3 ------------------%
\begin{homework}[3]
用傅里叶变换求解三维热传导方程的柯西问题
\[
\begin{cases}
\dfrac{\partial u}{\partial t}
=a^2\left(
\dfrac{\partial^2 u}{\partial x^2}
+\dfrac{\partial^2 u}{\partial y^2}
+\dfrac{\partial^2 u}{\partial z^2}
\right),
& t>0,\ (x,y,z)\in\mathbb R^3,\\[6pt]
u(x,y,z,0)=\varphi(x,y,z),
& (x,y,z)\in\mathbb R^3.
\end{cases}
\]
\end{homework}

\section{4极值原理——三题}
p93
\begin{homework}[2]
利用证明热传导方程极值原理的方法，证明：满足方程
\[
\frac{\partial^2 u}{\partial x^2}+\frac{\partial^2 u}{\partial y^2}=0
\]
的函数在有界闭区域上的最大值不超过它在边界上的最大值。
\end{homework}





\begin{homework}[3]
设 $u(x,t)$ 为热传导方程
\[
u_t - a^2 u_{xx} - c u = 0,\qquad c>0,
\]
在矩形区域
\[
R=\{(x,t)\mid 0 < x < l,\ 0<t<T\}
\]
中的解。

已知：
\[
|u(0,t)|,\ |u(l,t)| \le M,\qquad t\in[0,T],
\]
\[
|u(x,0)|\le M,\qquad x\in[0,l].
\]

试证：
\[
|u(x,t)|\le M e^{ct},\qquad (x,t)\in R.
\]
\end{homework}



\begin{homework}[5]
设 $u(x,t)$ 是区域 $\{0\le x\le \pi,\ 0\le t<\infty\}$ 中问题的经典解：
\[
\begin{cases}
u_t=u_{xx},\\[4pt]
u_x|_{x=0}=u_x|_{x=\pi}=0,\\[4pt]
u|_{t=0}=\varphi(x),
\end{cases}
\]
其中 $\varphi(0)=\varphi'(\pi)=0$。

证明：
\[
\sup_{0<x<\pi}|u(x,1)|\ \le\ \sup_{0<x<\pi}|\varphi(x)|.
\]

\textit{提示：将函数 $u(x,t)$ 关于直线 $x=\pi$ 作偶延拓到 $x\in(\pi,2\pi)$ 中。}
\end{homework}





   

\section{12.1 5解的渐进太六题}
95

\begin{homework}[2]
证明：当 $\varphi(x,y)$ 为 $\mathbb R^2$ 上的有界连续函数，且 $\varphi\in L^1(\mathbb R^2)$ 时，二维热传导方程柯西问题的解，当 $t\to\infty$ 时，以 $t^{-1}$ 衰减率趋于零。
\end{homework}




\begin{homework}[3]
证明：当 $\varphi(x,y,z)$ 为 $\mathbb R^3$ 上的有界连续函数，且 $\varphi\in L^1(\mathbb R^3)$ 时，三维热传导方程柯西问题的解，当 $t\to\infty$ 时，以 $t^{-3/2}$ 衰减率趋于零。
\end{homework}

105-106

\begin{homework}[1]
设
\[
u(x_1,x_2,\dots,x_n)=f(r),\qquad
r=\sqrt{x_1^2+x_2^2+\cdots+x_n^2},
\]
是 $n$ 维调和函数（即满足方程
\[
\frac{\partial^2 u}{\partial x_1^2}
+\frac{\partial^2 u}{\partial x_2^2}
+\cdots
+\frac{\partial^2 u}{\partial x_n^2}=0),
\]
试证明：
\[
f(r)=c_1+\frac{c_2}{r^{\,n-2}}\quad (n\ne 2),\qquad
f(r)=c_1+c_2\ln\frac1r\quad (n=2),
\]
其中 $c_1,c_2$ 为任意常数。
\end{homework}

\begin{dxtips}
 1.会求二次偏导
 2.会解常微分方程，常微分方程要用两次变量替换一个是降维把f（x）’=g（x）第二积分的时候变量替换把g（x)'挪到积分号后面。
\end{dxtips}
    





\begin{homework}[2]
证明：拉普拉斯算子在球面坐标 $(r,\theta,\varphi)$ 下可以写成
\[
\Delta u
=\frac1{r^2}\frac{\partial}{\partial r}\!\left(r^2\frac{\partial u}{\partial r}\right)
+\frac1{r^2\sin\theta}\frac{\partial}{\partial\theta}\!\left(\sin\theta\frac{\partial u}{\partial\theta}\right)
+\frac1{r^2\sin^2\theta}\frac{\partial^2 u}{\partial\varphi^2}.
\]
\end{homework}


\begin{homework}[3]
证明：拉普拉斯算子在柱面坐标 $(r,\theta,z)$ 下可以写成
\[
\Delta u
=\frac1r\frac{\partial}{\partial r}\!\left(r\frac{\partial u}{\partial r}\right)
+\frac1{r^2}\frac{\partial^2 u}{\partial\theta^2}
+\frac{\partial^2 u}{\partial z^2}.
\]
\end{homework}


\begin{homework}[6]
用分离变量法求解由下述调和方程的第一边值问题所描述的矩形平板
\[
(0\le x\le a,\;0\le y\le b)
\]
上的稳态温度分布：
\[
\begin{cases}
\dfrac{\partial^2 u}{\partial x^2}
+\dfrac{\partial^2 u}{\partial y^2}=0,\\[0.6em]
u(0,y)=u(a,y)=0,\\[0.4em]
u(x,0)=\sin\dfrac{\pi x}{a},\quad u(x,b)=0.
\end{cases}
\]
\end{homework}


\section{6格林公式四题}
\begin{homework}[1]
证明：(2.7) 式对于 $M_0$ 在 $\Omega$ 外与 $\Gamma$ 上的情形成立。
\end{homework}
\begin{equation}\tag{2.7}\label{eq:2.7}
 - \iint_{\Gamma}
 \left[
   u\,\frac{\partial}{\partial n}\!\left(\frac{1}{r}\right)
   - \frac{1}{r}\,\frac{\partial u}{\partial n}
 \right]\mathrm{d}S
 =
 \begin{cases}
   0, & \text{若 } M_0 \text{ 在 } \Omega \text{ 外},\\[4pt]
   2\pi u(M_0), & \text{若 } M_0 \text{ 在 } \Gamma \text{ 上},\\[4pt]
   4\pi u(M_0), & \text{若 } M_0 \text{ 在 } \Omega \text{ 内}.
 \end{cases}
\end{equation}



\begin{homework}[2]
若函数 $u(x,y)$ 是\textcolor{purple}{单位圆}上的调和函数，又已知它在\textcolor{purple}{单位圆周上}的数值为
\[
u = \sin\theta,
\]
其中 $\theta$ 表示极角，问函数 $u$ 在原点之值等于多少？
\end{homework}


\begin{homework}[3]
设 $u(M)$ 在 $\Omega$ 内调和，$M_0$ 是 $\Omega$ 中的任意一点。
$B_a$ 是以 $M_0$ 为球心、$a$ 为半径的球体，其体积为 $|B_a|$，证明：
\[
u(M_0)=\frac{1}{|B_a|}\iiint_{B_a} u\,dV.
\]
\end{homework}




\begin{homework}[4]
证明：当 $u(M)$ 在闭曲面 $\Gamma$ 的外部调和，并且在无穷远处成立
\[
u(M)=O\!\left(\frac{1}{r_{OM}}\right),\qquad
\frac{\partial u}{\partial r}=O\!\left(\frac{1}{r_{OM}^2}\right)
\quad (r_{OM}\to\infty),
\]
而 $M_0$ 是 $\Gamma$ 外的任一点，则公式 (2.6) 仍成立。
\begin{equation}\tag{2.6}\label{eq:2.6}
u(M_0)
= -\frac{1}{4\pi}
  \iint_{\Gamma}
  \left[
    u(M)\,\frac{\partial}{\partial n}
      \!\left(\frac{1}{r_{M_0M}}\right)
    - \frac{1}{r_{M_0M}}\,
      \frac{\partial u(M)}{\partial n}
  \right]\mathrm dS_M .
\end{equation}
\end{homework}



\section{7}
P121
\begin{homework}[1]
证明格林函数的性质 3 及性质 5。
\end{homework}

\begin{homework}[5]
求半圆区域上狄利克雷问题的格林函数。
\end{homework}